%======================================================================
%  GPN-Plasmids -- manuscript
%
%  Template derived from the GPN-Star manuscript (reference/GPN-Star/main.tex).
%  The preamble below is carried over verbatim so that the two manuscripts
%  share formatting, cross-reference macros and bibliography behaviour.
%
%  Everything marked  % TODO  is a placeholder for you to fill in.
%  See README.md for the conventions used in this file.
%======================================================================

\documentclass[11pt]{article}

\usepackage[utf8]{inputenc}
\usepackage{amsmath,amssymb}
\usepackage{amsfonts,amsthm}
\usepackage{bbm}
\usepackage[font=small]{caption}
\DeclareCaptionLabelFormat{previous-page}{#1 #2 \emph{(continued)}}
\DeclareUnicodeCharacter{2032}{\ensuremath{'}}  % ' prime  (U+2032), e.g. "5' UTR"
\DeclareUnicodeCharacter{2217}{\ensuremath{*}}  % * asterisk operator (U+2217)
\DeclareUnicodeCharacter{2212}{\ensuremath{-}}  % - minus sign (U+2212)

\usepackage{rotating}
\usepackage{graphicx}
\usepackage{xcolor}
\usepackage{url}
\usepackage[toc,page]{appendix} %for appendix/SI numbering

\usepackage{float} % for figures to stay where I want them to appear
% \usepackage[nolists, nomarkers]{endfloat}
\usepackage{bm} % for bold math symbols
\usepackage[margin=1in]{geometry}
\usepackage[backend=biber,style=nature,defernumbers=true]{biblatex}

\addbibresource{songlab.bib}

% Two bibliographies: everything cited in the main text (up to
% \maintextpartfalse) is collected in the "maincited" category and printed as
% "References"; anything cited only in the supplement is printed at the very
% end as "Supplementary References".
\DeclareBibliographyCategory{maincited}
\newif\ifmaintextpart\maintextparttrue
\AtEveryCitekey{\ifmaintextpart\addtocategory{maincited}{\thefield{entrykey}}\fi}

\usepackage[colorlinks=true,
            linkcolor=blue,
            urlcolor=blue,
            citecolor=blue]{hyperref}
\usepackage{xurl}
\usepackage{algorithm,algorithmic}
\usepackage{booktabs}
\usepackage{makecell}
\usepackage{multirow}


\theoremstyle{plain}
\newtheorem{axiom}{Axiom}
\newtheorem{claim}[axiom]{Claim}
\newtheorem{theorem}{Theorem}[section]
\newtheorem{lemma}[theorem]{Lemma}

\newcommand{\footremember}[2]{%
    \footnote{#2}
    \newcounter{#1}
    \setcounter{#1}{\value{footnote}}%
}
\newcommand{\footrecall}[1]{%
    \footnotemark[\value{#1}]%
}

% ---- cross-reference helpers -----------------------------------------
\newcommand{\tref}[1]{Table~\ref{#1}}
\newcommand{\fref}[1]{Figure~\ref{#1}}
\newcommand{\supptref}[1]{Supplementary Table~\ref{#1}}
\newcommand{\suppfref}[1]{Supplementary Figure~\ref{#1}}
\newcommand{\supptrefs}[1]{Supplementary Tables~\ref{#1}}
\newcommand{\suppfrefs}[2]{Supplementary Figures~\ref{#1} and \ref{#2}}
\newcommand{\edfref}[1]{Extended Data Figure~\ref{#1}}
\newcommand{\edfrefs}[1]{Extended Data Figures~\ref{#1}}
\newcommand{\edtref}[1]{Extended Data Table~\ref{#1}}

\newcommand{\red}{\color{red}}
\newcommand{\green}{\color{green}}
\newcommand{\blue}{\color{blue}}
\newcommand{\purple}{\color{purple}}

% ---- math shorthands --------------------------------------------------
\newcommand{\bsD}{\boldsymbol{D}}
\newcommand{\bsM}{\boldsymbol{M}}
\newcommand{\bstheta}{\boldsymbol{\theta}}
\newcommand{\bsbeta}{\boldsymbol{\beta}}
\newcommand{\bsgamma}{\boldsymbol{\gamma}}
\newcommand{\E}{\mathbb{E}}
\newcommand{\Var}{\text{Var}}
\newcommand{\Prob}{\mathbb{P}}
\newcommand{\bsc}{\boldsymbol{c}}
\newcommand{\bsC}{\boldsymbol{C}}
\newcommand{\bsS}{\boldsymbol{S}}
\newcommand{\bsw}{\boldsymbol{w}}
\newcommand{\bsx}{\boldsymbol{x}}
\newcommand{\bsv}{\boldsymbol{v}}
\newcommand{\bsu}{\boldsymbol{u}}
\newcommand{\bsy}{\boldsymbol{y}}
\newcommand{\bsz}{\boldsymbol{z}}
\newcommand{\bsY}{\boldsymbol{Y}}
\newcommand{\bsX}{\boldsymbol{X}}
\newcommand{\bsB}{\boldsymbol{B}}
\newcommand{\bsZ}{\boldsymbol{Z}}
\newcommand{\bsmu}{\boldsymbol{\mu}}
\newcommand{\bstau}{\boldsymbol{\tau}}

\def\bbP{\mathbb{P}}

% ---- float numbering for the supplement / extended data ---------------
\newcommand{\beginsupplement}{%\
        \renewcommand{\tablename}{Supplementary Table}
        \setcounter{table}{0}
        \renewcommand{\thetable}{\arabic{table}}%
        \renewcommand{\figurename}{Supplementary Figure}
        \setcounter{figure}{0}
        \renewcommand{\thefigure}{\arabic{figure}}%
        \setcounter{page}{1}
     }
\newcommand{\beginextendeddata}{%\
        \renewcommand{\tablename}{Extended Data Table}
        \setcounter{table}{0}
        \renewcommand{\thetable}{\arabic{table}}%
        \renewcommand{\figurename}{Extended Data Figure}
        \setcounter{figure}{0}
        \renewcommand{\thefigure}{\arabic{figure}}%
     }

\usepackage{xspace}
% Model / method name used throughout the manuscript. Change it in one place
% here and it updates everywhere.
\newcommand{\model}{GPN-Plasmids\xspace}   % TODO: confirm the final method name

\def\todo#1{\textbf{\textcolor{red}{ [TODO: #1]}}}
\def\fixme#1{\textbf{\textcolor{red}{ [FIXME: #1]}}}

% ---- placeholder graphic ---------------------------------------------
% Draws a labelled empty box so the skeleton compiles before any figure PDFs
% exist. Replace each \placeholderfig with \includegraphics and delete this
% definition once all figures are in figures/.
\newcommand{\placeholderfig}[3][\textwidth]{%
  \fbox{\parbox[c][#2][c]{\dimexpr#1-2\fboxsep-2\fboxrule\relax}{%
    \centering\ttfamily\small #3}}%
}

% ---- external documents (xr) -----------------------------------------
% Only needed if the supplement is split into a separate .tex/.pdf; kept from
% GPN-Star in case you go that route.
\usepackage{xr}

\makeatletter
\newcommand*{\addFileDependency}[1]{% argument=file name and extension
\typeout{(#1)}% latexmk will find this if $recorder=0
\@addtofilelist{#1}
\IfFileExists{#1}{}{\typeout{No file #1.}}
}\makeatother

\newcommand*{\myexternaldocument}[1]{%
\externaldocument{#1}%
\addFileDependency{#1.tex}%
\addFileDependency{#1.aux}%
}

%\myexternaldocument{supplementary}

\usepackage{nameref}

% Uncomment to stamp a watermark on drafts circulated outside the lab.
% \usepackage[printwatermark]{xwatermark}
% \newwatermark[pages=0-100,color=red!50,angle=0,scale=0.7,xpos=0,ypos=-125]{Confidential.  Please do not share.}


%======================================================================
%  Title and authors
%======================================================================

% TODO: settle on a title. Keep discarded candidates commented out here, as in
% GPN-Star, so the group can compare them.
% \title{A genomic language model for plasmids}
\title{GPN-Plasmids: \todo{title}}

% ---- author list ------------------------------------------------------
% Authors are HIDDEN: nothing is typeset below the title, and the PDF metadata
% carries no author either. The list is kept in the source so it is not lost.
% Flip to \showauthorstrue to print it.
\newif\ifshowauthors
\showauthorsfalse

\ifshowauthors
  % TODO: complete the author list, the equal-contribution footnote and the
  % affiliation numbering before this is ever switched on.
  \author{Jianan Canal Li$^{1,*}$, \todo{co-authors}, Yun S. Song$^{1,2,3,4,}$\footnote{To whom correspondence should be addressed: \url{yss@berkeley.edu}}\\
  \and
  \small$^{1}$Computer Science Division, University of California, Berkeley, Berkeley, CA, USA\and
  \small$^{2}$Department of Statistics, University of California, Berkeley, Berkeley, CA, USA\and
  \small$^{3}$Center for Computational Biology, University of California, Berkeley, Berkeley, CA, 94720, USA\\
  \small$^{4}$Innovative Genomics Institute, University of California, Berkeley, Berkeley, CA, 94720, USA\\
  }
\else
  \author{}
  \hypersetup{pdfauthor={}}
\fi

\date{}

\begin{document}

\maketitle


\begin{abstract}
% word limit 200 (Nature); GPN-Star ran ~220 words in draft
\todo{Abstract.}
% Suggested shape, following GPN-Star:
%  1. Context: what problem in plasmid biology / genome interpretation is open.
%  2. Gap: why existing methods fall short.
%  3. "Here, we introduce \model (...expansion of the acronym...), a ..."
%  4. What it was trained on.
%  5. Headline results, one clause per benchmark family.
%  6. Generalization beyond the primary setting.
%  7. "Taken together, these results position \model as ..."
\end{abstract}

\clearpage

%======================================================================
\section*{Introduction}
%======================================================================
% GPN-Star used five paragraphs here:
%  (1) the biological problem, (2) why language models are a natural fit and
%  where they fall short, (3) the alternative line of work being built on,
%  (4) "Here, we present \model ..." with the main claims and forward
%  references to figures, (5) a short paragraph on generality.

\todo{Introduction.}


%======================================================================
\section*{Results}
%======================================================================
% One \subsection* per result, each a short noun phrase (GPN-Star used
% "Pathogenic coding variants", "Complex trait heritability", ...). Each
% subsection opens by naming the question, then reports the comparison and
% points at a panel with \fref{...}.

\subsection*{\todo{Model and data overview}}

\todo{Describe the model, the training data, and the evaluation setup
(\fref{fig:overview}).}

\subsection*{\todo{Result 2}}

\todo{Text.}

\subsection*{\todo{Result 3}}

\todo{Text.}


%======================================================================
\section*{Discussion}
%======================================================================
% GPN-Star: summary of contributions, then limitations, then outlook on
% growing data resources.

\todo{Discussion.}


%======================================================================
\section*{Acknowledgements}
%======================================================================
\todo{Acknowledgements.}
% Remember to credit data providers and, if applicable, icon sources, e.g.:
% In \fref{fig:overview}, icons are made by ... from www.flaticon.com.

\section*{Funding}
% TODO: verify the grant numbers that apply to this project.
This research is supported in part by National Institutes of Health grants
\todo{grant numbers}.


\clearpage
\printbibliography[category=maincited,title={References}]
\newpage


%======================================================================
%  Main figures
%======================================================================
% Figures live at the end of the main text, after the references, as required
% by Nature-style submissions. Replace \placeholderfig with
% \includegraphics[width=\textwidth]{figures/<name>} once the PDF exists.
% Caption convention: a bold lead sentence naming the figure, then
% \textbf{(A)} ... one sentence per panel, then a closing sentence giving the
% metric and the error bars.

\begin{figure}[h]
    \centering
    \placeholderfig{6cm}{figures/overview}
    % \includegraphics[width=\textwidth]{figures/overview}
\caption{
\textbf{Overview of \model.} \textbf{(A)} \todo{panel A.} \textbf{(B)} \todo{panel B.} \textbf{(C)} \todo{panel C.}
}
\label{fig:overview}
\end{figure}

\begin{figure}[h]
    \centering
    \placeholderfig{6cm}{figures/benchmarks}
    % \includegraphics[width=\textwidth]{figures/benchmarks}
\caption{
\textbf{\todo{Benchmark figure title}.} \textbf{(A)} \todo{panel A.}
The performance metric is \todo{metric}. The error bars represent 95\% confidence intervals from 1,000 class-stratified bootstrap resamples.
}
\label{fig:benchmarks}
\end{figure}


\clearpage

%======================================================================
\section*{Methods}
\label{sec:methods}
%======================================================================
% In GPN-Star the Methods section is a one-paragraph pointer to the
% Supplementary Information, followed by the journal-required statements.
% Cross-reference it from the text with \nameref{sec:methods}.

The Supplementary Information contains Supplementary Text and the details of
\todo{list the supplementary methods subsections here}.

\subsection*{Reporting summary}
Further information on research design is available in the Nature Portfolio Reporting Summary linked to this article.

\subsection*{Data Availability}
\todo{Hugging Face / accession links for models, training data and benchmarks.}

\subsection*{Code Availability}
\todo{GitHub link and Zenodo archive DOI.}

\subsection*{Author Contributions}
\todo{Contribution statement, one clause per author, semicolon separated.}

\subsection*{Competing Interests}
The authors declare no competing interests.

\subsection*{Additional Information}

\noindent
\textbf{Supplementary information} The online version contains supplementary material.

\noindent
% Named here in GPN-Star, but a name in this line would leak the corresponding
% author while \showauthorsfalse is in effect. Fill in at submission.
\textbf{Correspondence and requests for materials} should be addressed to \todo{corresponding author}.


\clearpage

%======================================================================
\section*{Extended Data Figures}
%======================================================================
% Resets the figure counter and relabels floats as "Extended Data Figure N".

\beginextendeddata

\begin{figure}[h!]
    \centering
    \placeholderfig[0.9\textwidth]{6cm}{figures/ed-example}
    % \includegraphics[width=0.9\textwidth]{figures/ed-example}
\caption{
\todo{Extended data caption. These captions are not bold-led in GPN-Star; they open directly with the description.}
}
\label{fig:ed-example}
\end{figure}


\clearpage

%======================================================================
%  SUPPLEMENTARY INFORMATION
%======================================================================
% \maintextpartfalse must come before any supplement-only citation: from here
% on, newly cited references go to "Supplementary References" instead of the
% main "References" list.
\maintextpartfalse

% Resets figure/table counters and page numbering for the supplement.
\beginsupplement

\section*{Supplementary Text}
\label{sec:supp_text}
% Cross-reference from the main text with \nameref{sec:supp_text}.

\subsection*{\todo{Supplementary discussion topic}}

\todo{Text.}


\section*{Supplementary Methods}

\subsection*{Training Data}
\todo{Text.}

\subsection*{Model Architecture}
\todo{Text.}

\subsection*{Training Setup}
\todo{Text.}

\subsection*{Inference}
\todo{Text.}

\subsection*{Evaluation Benchmarks}
\todo{Text.}


\clearpage
\section*{Supplementary Tables}
\

% Wide tables use sidewaystable (rotating package) + \resizebox, as in the
% GPN-Star training-setup table.
\begin{table}[h]
\caption{
\todo{Table caption.}
}
\label{tab:example}
\begin{center}
\begin{tabular}{lcc}
\toprule
 & \todo{col 1} & \todo{col 2} \\
\midrule
\todo{row} & & \\
\bottomrule
\end{tabular}
\end{center}
\end{table}


\clearpage
\section*{Supplementary Figures}

\begin{figure}[h]
    \centering
    \placeholderfig[0.8\textwidth]{6cm}{figures/supp-example}
    % \includegraphics[width=0.8\textwidth]{figures/supp-example}
\caption{
\todo{Supplementary figure caption.}
}
\label{fig:supp-example}
\end{figure}


\clearpage
\printbibliography[notcategory=maincited,title={Supplementary References}]

\end{document}
